\documentclass[9pt,a4paper]{article}
\usepackage[paperwidth=210mm,paperheight=297mm,top=10mm,bottom=10mm,left=12mm,right=12mm]{geometry}
\usepackage{xcolor}
\usepackage{graphicx}
\usepackage{hyperref}
\usepackage{ragged2e}
\usepackage{fontspec}
\usepackage{fancyhdr}
\usepackage{lastpage}
\setmainfont[BoldFont=Lato-Bold,ItalicFont=Lato-Italic,BoldItalicFont=Lato-BoldItalic]{Lato-Regular}
\setsansfont[BoldFont=Lato-Bold,ItalicFont=Lato-Italic,BoldItalicFont=Lato-BoldItalic]{Lato-Regular}
\setlength{\parindent}{0pt}
\setlength{\parskip}{0pt}
\setlength{\fboxsep}{0pt}
\definecolor{textfg}{HTML}{20262D}
\definecolor{mutedfg}{HTML}{59646F}
\definecolor{accent}{HTML}{2D5F9D}
\definecolor{accentstrong}{HTML}{1F4476}
\definecolor{rulefg}{HTML}{D7DDE4}
\definecolor{surface}{HTML}{F4F6F9}
\definecolor{badgedark}{HTML}{2A2D31}
\definecolor{badgegreen}{HTML}{5B8648}
\definecolor{badgeblue}{HTML}{2D5F9D}
\definecolor{badgeslate}{HTML}{5F6C7D}
\color{textfg}
\hypersetup{colorlinks=true, urlcolor=accent}

% ── Page style ──────────────────────────────────────────────
\pagestyle{fancy}
\fancyhf{}
\renewcommand{\headrulewidth}{0pt}
\fancyfoot[L]{{\color{mutedfg}\fontsize{6}{7}\selectfont Toni Gansel --- Projekthistorie}}
\fancyfoot[R]{{\color{mutedfg}\fontsize{6}{7}\selectfont Seite \thepage{} von \pageref{LastPage}}}

% ── Macros ──────────────────────────────────────────────────
\newcommand{\phname}[1]{{\fontsize{16}{17}\selectfont\bfseries #1\par}}
\newcommand{\phheadline}[1]{{\color{accent}\fontsize{8.6}{9.2}\selectfont\bfseries #1\par}}
\newcommand{\phmetaline}[1]{{\color{mutedfg}\fontsize{7}{8}\selectfont #1\par}}
\newcommand{\phdoctitle}[1]{{\color{accent}\fontsize{9}{10}\selectfont\bfseries\MakeUppercase{#1}\par}}
\newcommand{\phprojectname}[1]{{\fontsize{8.5}{9.2}\selectfont\bfseries #1}}
\newcommand{\phperiod}[1]{{\color{accent}\fontsize{6.5}{7}\selectfont\bfseries\MakeUppercase{#1}}}
\newcommand{\phprojectmeta}[1]{{\color{mutedfg}\fontsize{7}{8}\selectfont #1\par}}
\newcommand{\phcontext}[1]{\vspace{0.8mm}{\setlength{\fboxsep}{1.5mm}\colorbox{surface}{\parbox{\dimexpr\linewidth-2\fboxsep-2pt\relax}{\color{mutedfg}\fontsize{7}{8.2}\selectfont #1}}}\par\vspace{0.8mm}}
\newcommand{\phresptitle}[1]{{\fontsize{7.5}{8.4}\selectfont\bfseries #1\par}}
\newcommand{\phrespdesc}[1]{{\color{mutedfg}\fontsize{7}{8.2}\selectfont #1\par}}
\newcommand{\phtoolline}[1]{{\color{accent}\fontsize{6.8}{7.6}\selectfont\bfseries #1\par}}

\begin{document}
\sloppy

% ── Header ──────────────────────────────────────────────────
\noindent
\begin{minipage}[t]{0.65\textwidth}
\vspace*{0pt}
\phname{Toni Gansel}
\phheadline{Testingenieur mit Fokus auf Sicherheit, Performance und AI-gestützte Automatisierung}
\vspace{0.5mm}
\phmetaline{Am Kirchbrunnen 18, 38446, Wolfsburg | +49 156 797 434 06 | toni.gansel@gmail.com}
{\color{accentstrong}\fontsize{6.8}{7.6}\selectfont
\href{https://www.linkedin.com/in/toni-gansel-814813245/}{LinkedIn} \textbar{} \href{https://github.com/tonileet}{GitHub}\par}
\end{minipage}
\hfill
\begin{minipage}[t]{0.3\textwidth}
\vspace*{0pt}
\raggedleft
\phdoctitle{Projekthistorie}
{\color{mutedfg}\fontsize{7}{8}\selectfont 24 Projekte\par}
\end{minipage}

\vspace{2mm}
{\color{accent}\rule{\linewidth}{0.6pt}\par}
\vspace{3mm}

% ── Projects ────────────────────────────────────────────────

% Project heading
\noindent
{\phprojectname{Elektronische Patientenakte (EPA)}\hfill\phperiod{01/2024 - 01/2026}\par}
{\phprojectmeta{RISE | Performance Tester | Continuous Integration und Agile Vorgehensweise}%
\hfill {\setlength{\fboxsep}{1.2pt}\colorbox{badgeblue}{\textcolor{white}{\fontsize{5.5}{6.0}\selectfont\bfseries AI}}}
\par}

\phcontext{Im Rahmen der Einführung der Elektronischen Patientenakte (EPA) war ich als Performance Tester verantwortlich für die Planung, Entwicklung und Auswertung umfangreicher Last- und Performancetests. Die EPA bildet die zentrale Plattform zur Verwaltung und Speicherung der Gesundheitsdaten aller deutschen Versicherten und stellt damit höchste Anforderungen an Skalierbarkeit, Verfügbarkeit und Stabilität. Zur Testdurchführung wurde ein JMeter-basiertes Framework eingesetzt, das über Groovy auf eine selbstentwickelte Java-Bibliothek zugreift. Die Tests wurden in einem Kubernetes-Cluster ausgeführt und mit Prometheus sowie Grafana überwacht. Ergänzend wurden KI-gestützte Methoden zur Programmierunterstützung und im Testmanagement erprobt sowie automatisierte Lasttestberichte mit R erstellt.}

\vspace{0.5mm}
\phresptitle{Planung der Entwicklung eines Lasttests mit Jira}
\phrespdesc{Ich plante die Entwicklung der Lasttests strukturiert mit Jira, erstellte einen detaillierten Entwicklungsplan, definierte Testziele und Abnahmekriterien und koordinierte die Abstimmung der Lastszenarien und Testumgebungen mit allen beteiligten Teams.}
\vspace{0.8mm}
\phresptitle{Entwicklung eines JMeter Testplans mit Zugriff auf eine selbstentwickelte Java Bibliothek via Groovy}
\phrespdesc{Ich entwarf und implementierte einen komplexen JMeter-Testplan, der über Groovy-Scripting auf eine intern entwickelte Java-Bibliothek zugreift. So integrierte ich projektspezifische Authentifizierungsmechanismen und Datengenerierungslogiken nahtlos in die Lasttests und erhöhte die Testabdeckung und Realitätsnähe der Szenarien deutlich.}
\vspace{0.8mm}
\phresptitle{Testausführung aus einem Kubernetes-Cluster}
\phrespdesc{Ich konfigurierte, orchestrierte und führte die Lasttests in einer produktionsnahen Kubernetes-Umgebung durch. Die verteilte Ausführung ermöglichte mir, hohe Nutzerzahlen unter realistischen Bedingungen zu simulieren. Das Monitoring der Systemmetriken führte ich in Echtzeit über Prometheus und Grafana durch, um Engpässe frühzeitig zu erkennen.}
\vspace{0.8mm}
\phresptitle{Erforschung von KI-gestützter Programmierung und Testmanagement}
\phrespdesc{Ich analysierte und erprobte systematisch KI-gestützte Werkzeuge zur Unterstützung bei Testplanung, Codegenerierung und Auswertung von Testergebnissen. Ich evaluierte den Mehrwert und die Grenzen aktueller KI-Assistenten im Testkontext und erarbeitete Empfehlungen für deren Einsatz im Projektalltag.}
\vspace{0.8mm}
\phresptitle{Automatisierte Erstellung von Lasttestberichten mit R}
\phrespdesc{Ich entwickelte R-Skripte zur automatisierten Aufbereitung, Auswertung und Visualisierung von Lasttestdaten. Die von mir generierten Berichte ermöglichten eine schnelle und konsistente Kommunikation der Testergebnisse an Stakeholder und reduzierten den manuellen Aufwand bei der Berichterstattung erheblich.}
\vspace{0.8mm}

\phtoolline{JMeter, Groovy, Kubernetes, GitLab, Prometheus, Grafana, Java, Jira, Docker, R, Java, Git, JetBrains IntelliJ}

\vspace{1.5mm}
{\color{rulefg}\rule{\linewidth}{0.25pt}\par}
\vspace{2.5mm}


% Project heading
\noindent
{\phprojectname{EMMA}\hfill\phperiod{06/2024 - 07/2024}\par}
{\phprojectmeta{RISE | Performance Tester | SCRUM}%
\hfill {\setlength{\fboxsep}{1.2pt}\colorbox{badgeblue}{\textcolor{white}{\fontsize{5.5}{6.0}\selectfont\bfseries AI}}}
\par}

\phcontext{EMMA – E-Mail Manager ist eine Softwarelösung zur zentralen Verwaltung und Verarbeitung geschäftlicher E-Mail-Kommunikation. Die Plattform unterstützt Unternehmen bei der strukturierten Organisation, Verarbeitung und Archivierung von E-Mails sowie bei der Integration in bestehende Geschäftsprozesse und IT-Systeme. Ziel ist eine sichere, transparente und effiziente Kommunikation in komplexen IT-Umgebungen.}

\vspace{0.5mm}
\phresptitle{Planung und Auswertung von Performance- und Lasttests}
\phrespdesc{Ich plante eigenverantwortlich die Teststrategie für Performance- und Lasttests, definierte Lastszenarien und Akzeptanzkriterien und stimmte mich mit dem Entwicklungsteam ab. Nach der Testdurchführung wertete ich die Messergebnisse strukturiert aus und dokumentierte Auffälligkeiten sowie Handlungsempfehlungen.}
\vspace{0.8mm}
\phresptitle{Aufbau und Ausführung eines JMeter-Testplans}
\phrespdesc{Ich entwarf und implementierte einen vollständigen JMeter-Testplan, der die relevanten Nutzerszenarien des EMMA-Systems abbildet. Ich testete verschiedene Lastprofile und analysierte die Systemreaktion unter unterschiedlichen Belastungsstufen, um Engpässe frühzeitig zu identifizieren. Die Testergebnisse wurden via Jira direkt an das Entwicklungsteam kommuniziert.}
\vspace{0.8mm}

\phtoolline{JMeter, Jira, Git}

\vspace{1.5mm}
{\color{rulefg}\rule{\linewidth}{0.25pt}\par}
\vspace{2.5mm}


% Project heading
\noindent
{\phprojectname{TMS}\hfill\phperiod{04/2024 - 07/2024}\par}
{\phprojectmeta{RISE | Test Manager | SCRUM}%
\hfill {\setlength{\fboxsep}{1.2pt}\colorbox{badgeblue}{\textcolor{white}{\fontsize{5.5}{6.0}\selectfont\bfseries AI}}}
\par}

\phcontext{TMS (Transaction Management Service) ist eine zentrale API-Plattform zur Integration und Orchestrierung von Backend-Systemen. Der Service nimmt Requests von Clients entgegen, kapselt die Komplexität verschiedener Endsysteme und leitet Anfragen über standardisierte Schnittstellen an die jeweiligen Zielsysteme weiter. Dadurch wird die Integration vereinfacht, Systemkopplung reduziert und eine skalierbare, konsistente Kommunikation zwischen Anwendungen ermöglicht.}

\vspace{0.5mm}
\phresptitle{Erstellung eines Testkonzepts für eine Telematikinfrastruktur-API mit Postman und Bruno}
\phrespdesc{Ich konzipierte und erstellte ein vollständiges Testkonzept für eine REST-API der Telematikinfrastruktur. Ich definierte Testumfänge, Teststufen, Testmethoden und Verantwortlichkeiten und legte Qualitätsziele sowie Abnahmekriterien in Abstimmung mit Projektleitung und Entwicklungsteam fest.}
\vspace{0.8mm}
\phresptitle{Formulierung und Verwaltung von API-Testfällen in Zephyr Scale und Jira}
\phrespdesc{Ich erstellte detaillierte, nachvollziehbare Testfälle in Jira und Zephyr Scale auf Basis der API-Spezifikation. Dabei deckte ich sowohl Positiv- als auch Negativszenarien ab und formulierte die Testfälle so, dass verschiedene Teammitglieder sie reproduzierbar durchführen konnten.}
\vspace{0.8mm}
\phresptitle{Implementierung von Testfällen mit Postman und Bruno}
\phrespdesc{Ich implementierte die definierten Testfälle als automatisierte API-Tests mit Postman und Bruno. Ich entwickelte Testskripte zur Validierung von Anfragen, Antworten, Statuscodes und Datenintegrität und bereitete Testumgebungen sowie Testdaten für eine reproduzierbare Testdurchführung vor.}
\vspace{0.8mm}

\phtoolline{Postman, Bruno, Jira, Zephyr Scale, Git, Maven, JetBrains IntelliJ}

\vspace{1.5mm}
{\color{rulefg}\rule{\linewidth}{0.25pt}\par}
\vspace{2.5mm}


% Project heading
\noindent
{\phprojectname{TI-Portal}\hfill\phperiod{02/2024 - 05/2024}\par}
{\phprojectmeta{RISE | Test Automation Engineer | Agile Vorgehensweise, Lehrgespräch und SCRUM}%
\hfill {\setlength{\fboxsep}{1.2pt}\colorbox{badgeblue}{\textcolor{white}{\fontsize{5.5}{6.0}\selectfont\bfseries AI}}}
\par}

\phcontext{TI-Portal ist eine E-Commerce-Plattform innerhalb der RISE-Telematikinfrastruktur, auf der verschiedene Anbieter ihre Produkte und Services anbieten können. Jeder Anbieter betreibt einen eigenen Shop, während Kunden über eine zentrale Plattform einkaufen. Der Zugang erfolgt über das TI-Gateway, wodurch eine sichere Integration in die TI-Infrastruktur gewährleistet wird und ein skalierbarer digitaler Marktplatz entsteht.}

\vspace{0.5mm}
\phresptitle{Testautomatisierung eines Webshops mit Zephyr Scale und Cucumber}
\phrespdesc{Ich konzipierte und implementierte eine Testautomatisierungslösung für den TI-Portal-Webshop auf Basis des Keyword-Driven-Ansatzes mit Cucumber. Ich verwaltete die Testfälle in Zephyr Scale und verknüpfte sie mit dem automatisierten Framework Cypress, um eine durchgehende Nachverfolgbarkeit zwischen Testfall und Testergebnis zu gewährleisten.}
\vspace{0.8mm}
\phresptitle{Aufbau einer automatisierten Testpipeline mit GitLab CI/CD}
\phrespdesc{Ich entwarf und implementierte eine vollständig automatisierte Testpipeline in GitLab CI/CD, die bei jedem Commit automatisch die Testautomatisierungs-Suite ausführt. Die Pipeline umfasste Build-, Test- und Reporting-Stufen und ermöglichte dem Entwicklungsteam eine frühzeitige Rückmeldung über den Qualitätsstatus des Webshops.}
\vspace{0.8mm}
\phresptitle{Migration der bestehenden Selenium-Testsuite zu Cypress}
\phrespdesc{Ich analysierte die vorhandene Selenium-basierte Testautomatisierung und migrierte die Testfälle schrittweise zu Cypress. Dabei modernisierte ich die Tests, verbesserte die Wartbarkeit und erhöhte die Ausführungsgeschwindigkeit. Die Migration führte ich iterativ durch, um den laufenden Betrieb nicht zu unterbrechen.}
\vspace{0.8mm}
\phresptitle{Integration der Tests in Testpipeline und Zephyr Scale}
\phrespdesc{Ich integrierte die automatisierten Tests in die CI/CD-Testpipeline und band sie bidirektional an Zephyr Scale an, sodass Testergebnisse automatisch in das Testmanagement-Tool zurückgeschrieben wurden. Damit schuf ich vollständige Transparenz über den Teststatus aller automatisierten Testfälle.}
\vspace{0.8mm}
\phresptitle{Mentoring und Einarbeitung mehrerer Junior Tester}
\phrespdesc{Ich begleitete mehrere Junior Tester strukturiert und arbeitete sie in die Methoden und Werkzeuge der Testautomatisierung ein. Dazu führte ich Pair-Programming-Sessions und Code-Reviews durch, erstellte Dokumentation und gab regelmäßiges Feedback zur Förderung des fachlichen Wachstums der Teammitglieder.}
\vspace{0.8mm}

\phtoolline{Selenium, Cypress, Cucumber, GitLab, Zephyr Scale, Jira, TypeScript, Git, Maven, JetBrains IntelliJ}

\vspace{1.5mm}
{\color{rulefg}\rule{\linewidth}{0.25pt}\par}
\vspace{2.5mm}


% Project heading
\noindent
{\phprojectname{HSK API Tests}\hfill\phperiod{05/2023 - 09/2023}\par}
{\phprojectmeta{Secunet | Test Manager | Agile Vorgehensweise}%
\hfill {\setlength{\fboxsep}{1.2pt}\colorbox{badgeblue}{\textcolor{white}{\fontsize{5.5}{6.0}\selectfont\bfseries AI}}}
\par}

\phcontext{secunet Highspeedkonnektor (HSK) ist eine skalierbare, serverbasierte Zugangskomponente zur deutschen Telematikinfrastruktur (TI), die zentral in Rechenzentren betrieben wird und den sicheren Zugriff auf TI-Fachdienste ermöglicht. Der HSK ersetzt viele dezentrale Konnektoren und bildet die technische Basis des TI-Gateways. Im Projekt wurden die Schnittstellen und APIs des HSK getestet, um eine sichere und zuverlässige Kommunikation mit der TI sicherzustellen.}

\vspace{0.5mm}
\phresptitle{Automatisierte Tests einer REST-API mit SoapUI}
\phrespdesc{Ich plante und implementierte eigenverantwortlich automatisierte API-Tests für die HSK-REST-API mit SoapUI. Ich analysierte die API-Spezifikation, leitete relevante Testszenarien ab, erstellte Testsuiten und führte die Tests durch und wertete die Ergebnisse mit anschließender Dokumentation der Befunde in Jira aus.}
\vspace{0.8mm}
\phresptitle{Aufbau und Integration einer Jenkins-Testpipeline}
\phrespdesc{Ich konfigurierte und implementierte eine Jenkins-basierte CI-Testpipeline zur automatisierten und regelmäßigen Ausführung der API-Tests. Ich stellte sicher, dass bei jedem Build-Prozess alle definierten Testfälle ausgeführt und Testergebnisse automatisch dokumentiert und an das Team kommuniziert wurden.}
\vspace{0.8mm}

\phtoolline{SoapUI, Jenkins, Jira, Git}

\vspace{1.5mm}
{\color{rulefg}\rule{\linewidth}{0.25pt}\par}
\vspace{2.5mm}


% Project heading
\noindent
{\phprojectname{Cell Broadcast}\hfill\phperiod{04/2022 - 10/2022}\par}
{\phprojectmeta{Mecom | Test Manager | SCRUM}%
\hfill {\setlength{\fboxsep}{1.2pt}\colorbox{badgeslate}{\textcolor{white}{\fontsize{5.5}{6.0}\selectfont\bfseries Performance}}}
\par}

\phcontext{Cell Broadcast ist ein bundesweit eingesetztes Warnsystem zur schnellen Information der Bevölkerung bei Katastrophen und Gefahrenlagen. Warnmeldungen werden standortbasiert über Mobilfunkzellen direkt an alle kompatiblen Smartphones in einem betroffenen Gebiet gesendet – ohne App oder Registrierung. Die Lösung ermöglicht eine schnelle, skalierbare und netzlastunabhängige Verbreitung offizieller Warnungen.}

\vspace{0.5mm}
\phresptitle{Validierung von XML-Transformationen mit JUnit auf großen Datenmengen}
\phrespdesc{Ich entwickelte und führte automatisierte JUnit-Tests zur Validierung der XML-Transformationslogik der Kafka-Komponente durch. Dabei nutzte ich über 20.000 historische Nachrichten als Testbasis, um sicherzustellen, dass alle realen Nachrichtenformate korrekt transformiert und vollständig übernommen werden.}
\vspace{0.8mm}
\phresptitle{Testfälle mit fehlenden, manipulierten und typfremden Werten}
\phrespdesc{Ich entwickelte automatisierte Testfälle mit fehlenden, manipulierten und typfremden Werten, um die Robustheit der Anwendung zu prüfen. Dafür nutzte ich JUnit zur automatisierten Ausführung der Tests und verwendete XML-basierte Eingangs- und Ausgangsnachrichten, die ich gezielt veränderte, um Validierung, Fehlerbehandlung und die Stabilität der Schnittstellen zu überprüfen.}
\vspace{0.8mm}
\phresptitle{Sicherstellung vollständiger Datenübernahme bei Nachrichtentransformation}
\phrespdesc{Ich verifizierte systematisch mit JUnit, dass sämtliche relevanten Datenfelder bei der XML-Transformation vollständig und verlustfrei übernommen werden. Dazu führte ich strukturierte Vergleichstests zwischen Eingangs- und Ausgangsnachrichten durch, dokumentierte Abweichungen in Jira und kommunizierte sie an das Entwicklungsteam.}
\vspace{0.8mm}

\phtoolline{Java, JUnit, Jira, JetBrains IntelliJ, Kafka, XML, Git}

\vspace{1.5mm}
{\color{rulefg}\rule{\linewidth}{0.25pt}\par}
\vspace{2.5mm}


% Project heading
\noindent
{\phprojectname{school-SH}\hfill\phperiod{07/2019 - 06/2022}\par}
{\phprojectmeta{Dataport | Security Tester / Test Manager | SCRUM, Continuous Integration und Exploratives Testing}%
\hfill {\setlength{\fboxsep}{1.2pt}\colorbox{badgeblue}{\textcolor{white}{\fontsize{5.5}{6.0}\selectfont\bfseries AI}}}
\par}

\phcontext{School-SH ist eine landesweit eingesetzte Schulverwaltungssoftware des Landes Schleswig-Holstein zur Verwaltung zentraler schulischer Daten. Die Webanwendung unterstützt Schulen bei der Organisation von Schülerinnen und Schülern, Klassen, Kursen, Zeugnissen und statistischen Meldungen. Sie wird im Landesnetz betrieben und dient als zentrale Plattform zur digitalen Abbildung administrativer Schulprozesse.}

\vspace{0.5mm}
\phresptitle{Konzeption und Implementierung einer Teststrategie}
\phrespdesc{Ich erarbeitete und implementierte eine projektweite Teststrategie, die alle relevanten Testarten abdeckt. Ich definierte Testumfänge, Testmethoden, Verantwortlichkeiten und Qualitätsziele und legte damit das Fundament für alle weiteren Testaktivitäten im Projektverlauf.}
\vspace{0.8mm}
\phresptitle{Sicherheits-, Last- und Regressionstests mit OWASP ZAP, JMeter und Selenium}
\phrespdesc{Ich plante und führte umfassende Sicherheitstests mit OWASP ZAP nach der OWASP-Top-10-Methodik durch, simulierte mit JMeter realistische Lastszenarien und implementierte mit Selenium und Cucumber automatisierte, wartbare Regressionstests für die Schulverwaltungsplattform.}
\vspace{0.8mm}
\phresptitle{Projektbegleitende Schulungen zu Security und Testing}
\phrespdesc{Ich plante und führte projektbegleitende Schulungen für Entwickler und Tester zu den Themen Application Security und Softwaretesting durch. Dabei sensibilisierte ich das Team für sicherheitsrelevante Aspekte der Softwareentwicklung und stärkte die interne Testkompetenz für eine eigenständige Qualitätssicherung.}
\vspace{0.8mm}
\phresptitle{Aufbau einer automatisierten Testpipeline in Jenkins}
\phrespdesc{Ich konfigurierte und baute eine Jenkins-basierte CI/CD-Testpipeline auf, die die Regressionstests automatisiert ausführt. Die Pipeline integrierte ich in den bestehenden Entwicklungsprozess und ermöglichte damit eine kontinuierliche Qualitätskontrolle nach jedem Build, wodurch Regressionen frühzeitig erkannt wurden.}
\vspace{0.8mm}

\phtoolline{Jira, Jenkins, SonarQube, Tomcat, Java EE, PostgreSQL, OWASP ZAP, Selenium, Cucumber, JMeter, Nmap, Git, ANT, Eclipse}

\vspace{1.5mm}
{\color{rulefg}\rule{\linewidth}{0.25pt}\par}
\vspace{2.5mm}


% Project heading
\noindent
{\phprojectname{Login, Checkout}\hfill\phperiod{02/2019 - 04/2019}\par}
{\phprojectmeta{Fielmann | Test Manager | SCRUM}%
\hfill {\setlength{\fboxsep}{1.2pt}\colorbox{badgedark}{\textcolor{white}{\fontsize{5.5}{6.0}\selectfont\bfseries Security}}}
\par}

\phcontext{Im Rahmen der Weiterentwicklung des Fielmann-Online-Shops wurde eine Keycloak-basierte Identity-Lösung für Benutzer-Authentifizierung implementiert. Das Projekt umfasste Login-, Logout- und Passwort-Management-Funktionen. Anschliessend entstand eine mobile Anwendung zur virtuellen Anprobe von Sonnenbrillen, bei der Nutzer Modelle digital testen und anschließend direkt in der Appp kaufen können. Im Projekt formulierte ich die Teststrategie für den Checkout-Prozess.}

\vspace{0.5mm}
\phresptitle{Weiterentwicklung und Durchführung von E2E-Tests mit TestCafe}
\phrespdesc{Ich plante und implementierte eigenverantwortlich End-to-End-Tests für die Login- und Checkout-Funktionalitäten der Fielmann E-Commerce-Plattform mit TestCafe. Die Tests führte ich unter Berücksichtigung verschiedener Browser- und Gerätekonfigurationen durch und deckte damit alle kritischen Nutzerpfade vollständig ab.}
\vspace{0.8mm}
\phresptitle{Risikoanalyse mit OWASP Cornucopia}
\phrespdesc{Ich führte eine strukturierte Bedrohungsmodellierung und Risikoanalyse auf Basis des OWASP-Cornucopia-Kartenspiels durch. In gemeinsamen Workshop-Sessions mit dem Entwicklungsteam identifizierte, bewertete und priorisierte ich potenzielle Sicherheitsrisiken für die Authentifizierungs- und Checkout-Prozesse und definierte geeignete Gegenmaßnahmen.}
\vspace{0.8mm}
\phresptitle{Konzeption und Implementierung einer Teststrategie}
\phrespdesc{Ich erarbeitete eine umfassende Teststrategie für die zu prüfenden Systemkomponenten und legte Testziele, Testumfänge, einzusetzende Werkzeuge und Verantwortlichkeiten fest. Die Strategie stimmte ich mit der Projektleitung ab und schuf damit die Grundlage für eine strukturierte Qualitätssicherung des E-Commerce-Systems.}
\vspace{0.8mm}

\phtoolline{TestCafe, Docker, Keycloak, GIT, Jira, OWASP Cornucopia, JetBrains IntelliJ}

\vspace{1.5mm}
{\color{rulefg}\rule{\linewidth}{0.25pt}\par}
\vspace{2.5mm}


% Project heading
\noindent
{\phprojectname{DiViS}\hfill\phperiod{10/2017 - 12/2018}\par}
{\phprojectmeta{Hamburger Schulbehörde | Security Tester / Test Manager | SCRUM}%
\hfill {\setlength{\fboxsep}{1.2pt}\colorbox{badgedark}{\textcolor{white}{\fontsize{5.5}{6.0}\selectfont\bfseries Security}}}
\par}

\phcontext{DiViS (Digitale Verwaltung in Schulen) ist eine webbasierte Schulverwaltungssoftware der Freien und Hansestadt Hamburg zur digitalen Abbildung administrativer Schulprozesse. Die Anwendung unterstützt Schulen bei der Verwaltung von Schüler- und Klassendaten, Noten, Zeugnissen sowie organisatorischen Abläufen und stellt eine zentrale Datenbasis für Verwaltung und Bildungspolitik bereit.}

\vspace{0.5mm}
\phresptitle{Planung und Durchführung von Sicherheits- und Lasttests}
\phrespdesc{Ich konzipierte und führte Sicherheitstests auf Basis dem BSI Grundschutzkatalog mit OWASP ZAP durch und simulierte mit JMeter realistische Lastszenarien, um die Systemstabilität unter hoher Nutzerlast zu überprüfen. Dabei identifizierte ich Schwachstellen und Performanceengpässe und dokumentierte die Befunde nachvollziehbar in Testberichten und Jira.}
\vspace{0.8mm}
\phresptitle{Beratung der Projektleitung in Qualitäts- und Testfragen}
\phrespdesc{Ich beriet die Projektleitung fachlich in allen Fragen der Qualitätssicherung und der Teststrategie. Ich bewertete Risiken, gab Empfehlungen zu Testprioritäten, erstellte Statusberichte zum Testfortschritt und unterstützte bei Entscheidungen zu Testumfang und Ressourceneinsatz für eine effiziente Qualitätssicherung.}
\vspace{0.8mm}
\phresptitle{Weiterentwicklung der automatisierten Regressionstests}
\phrespdesc{Ich analysierte, pflegte und erweiterte die bestehende automatisierte Regressionstestsuite auf Basis von Ranorex mit C\#-Eigententwicklungen. Ich passte die Tests an neue Systemfunktionalitäten an, baute technische Schulden ab und erhöhte durch Refactoring und verbesserte Strukturierung die Wartbarkeit der Testskripte nachhaltig.}
\vspace{0.8mm}

\phtoolline{Jira, Jenkins, SonarQube, Tomcat, Java EE, PostgreSQL, OWASP ZAP, Ranorex, JMeter, Nmap, Git, ANT, Eclipse}

\vspace{1.5mm}
{\color{rulefg}\rule{\linewidth}{0.25pt}\par}
\vspace{2.5mm}


% Project heading
\noindent
{\phprojectname{DutyPay}\hfill\phperiod{04/2018 - 06/2018}\par}
{\phprojectmeta{TaxHub | Test Manager}%
\hfill {\setlength{\fboxsep}{1.2pt}\colorbox{badgegreen}{\textcolor{white}{\fontsize{5.5}{6.0}\selectfont\bfseries Automation}}}
\par}

\phcontext{Softwarelösung zur automatisierten Berechnung von Mehrwertsteuerverpflichtungen in E-Commerce-Systemen.}

\vspace{0.5mm}
\phresptitle{Aufbau eines Testkonzepts}
\vspace{0.8mm}
\phresptitle{Implementierung qualitätssichernder Maßnahmen}
\vspace{0.8mm}

\phtoolline{Jenkins, SonarQube, Apache HTTP Server, TYPO3}

\vspace{1.5mm}
{\color{rulefg}\rule{\linewidth}{0.25pt}\par}
\vspace{2.5mm}


% Project heading
\noindent
{\phprojectname{BusinessUnit Security}\hfill\phperiod{07/2017 - 02/2018}\par}
{\phprojectmeta{Sogeti | Test Manager}%
\hfill {\setlength{\fboxsep}{1.2pt}\colorbox{badgedark}{\textcolor{white}{\fontsize{5.5}{6.0}\selectfont\bfseries Security}}}
\par}

\phcontext{Erweiterung des Dienstleistungsportfolios um professionelle Security-Test-Services.}

\vspace{0.5mm}
\phresptitle{Aufbau der BusinessUnit Security}
\vspace{0.8mm}
\phresptitle{Koordination von Kundenprojekten im Bereich Security-Testing}
\vspace{0.8mm}
\phresptitle{Teamentwicklung und technische Leitung}
\vspace{0.8mm}

\phtoolline{Jira, Confluence}

\vspace{1.5mm}
{\color{rulefg}\rule{\linewidth}{0.25pt}\par}
\vspace{2.5mm}


% Project heading
\noindent
{\phprojectname{HSMS}\hfill\phperiod{06/2015 - 07/2017}\par}
{\phprojectmeta{Hamburger Schulbehörde | Test Manager | SCRUM}%
\hfill {\setlength{\fboxsep}{1.2pt}\colorbox{badgeslate}{\textcolor{white}{\fontsize{5.5}{6.0}\selectfont\bfseries Performance}}}
\par}

\phcontext{HSMS (Hamburger Schulverwaltungs Software) ist eine webbasierte Schulverwaltungssoftware der Freien und Hansestadt Hamburg zur digitalen Abbildung administrativer Schulprozesse. Die Anwendung unterstützt Schulen bei der Verwaltung von Schüler- und Klassendaten, Noten, Zeugnissen sowie organisatorischen Abläufen und stellt eine zentrale Datenbasis für Verwaltung und Bildungspolitik bereit.}

\vspace{0.5mm}
\phresptitle{Konzeption und Durchführung einer projektweiten Teststrategie}
\phrespdesc{Ich erarbeitete und entwickelte eine projektweite Teststrategie kontinuierlich weiter, die alle relevanten Testarten und Testphasen abdeckt. Ich definierte Testziele, Testumfänge, einzusetzende Werkzeuge und Verantwortlichkeiten und passte die Strategie regelmäßig an geänderte Anforderungen und Projektprioritäten an.}
\vspace{0.8mm}
\phresptitle{Planung und Administration der Testinfrastruktur}
\phrespdesc{Ich baute, konfigurierte und administrierte laufend die vollständige Testinfrastruktur aus Jenkins, JavaMelody, SonarQube, JMeter, OWASP ZAP und Ranorex. Ich stellte die dauerhafte Verfügbarkeit und Aktualität aller Testumgebungen und Werkzeuge sicher und führte notwendige Updates und Anpassungen eigenverantwortlich durch.}
\vspace{0.8mm}
\phresptitle{Sicherheits-, Regressions-, Last- und Performancetests sowie Codeanalysen}
\phrespdesc{Ich führte eigenständig alle im Rahmen der Teststrategie definierten Testaktivitäten durch: Sicherheitstests mit OWASP ZAP, automatisierte Regressionstests mit Ranorex, Lasttests mit JMeter und Codequalitätsanalysen mit SonarQube sowie manuelle Tests auf Basis definierter Testfälle.}
\vspace{0.8mm}
\phresptitle{Planung und Koordination regelmäßiger Endanwender-Tests}
\phrespdesc{Ich organisierte und koordinierte regelmäßige Usability- und Abnahmetests mit Endanwendern aus dem schulischen Umfeld. Ich plante die Test-Sessions, bereitete Testszenarien in Word vor, moderierte die Durchführung und bereitete die Ergebnisse strukturiert für Projektleitung und Entwicklungsteam auf.}
\vspace{0.8mm}

\phtoolline{Jira, Jenkins, SonarQube, JMeter, JavaMelody, OWASP ZAP, Ranorex, Eclipse, Git}

\vspace{1.5mm}
{\color{rulefg}\rule{\linewidth}{0.25pt}\par}
\vspace{2.5mm}


% Project heading
\noindent
{\phprojectname{SAST Integration}\hfill\phperiod{01/2015 - 05/2015}\par}
{\phprojectmeta{Volkswagen | Security Consultant | inkrementelles Modell}%
\hfill {\setlength{\fboxsep}{1.2pt}\colorbox{badgedark}{\textcolor{white}{\fontsize{5.5}{6.0}\selectfont\bfseries Security}}}
\par}

\phcontext{SAST Integration (Volkswagen Konzern) Konzeption und Implementierung eines Prototyps für automatisierte statische Quellcodeanalyse auf Basis von PMD und SonarQube. Zusätzlich wurde ein konzernweit standardisierter Regelkatalog entwickelt und mit diversen Stakeholdern abgestimmt, um Coding-Standards sowie sicherheitsrelevante Muster konsistent zu prüfen.}

\vspace{0.5mm}
\phresptitle{Implementierung eines SAST-Prototyps mit PMD und SonarQube}
\phrespdesc{Ich entwarf und implementierte einen funktionsfähigen Prototyp für die automatisierte statische Quellcodeanalyse auf Basis von PMD und SonarQube. Ich integrierte den Prototyp in bestehende Build-Prozesse und demonstrierte damit die technische Machbarkeit einer konzernweiten SAST-Lösung für den Volkswagen-Konzern.}
\vspace{0.8mm}
\phresptitle{Entwicklung eines konzernweiten Standardregelsatzes für PMD und SonarQube}
\phrespdesc{Ich analysierte bestehende Coding-Standards und Sicherheitsanforderungen im Volkswagen-Konzern und entwickelte daraus einen standardisierten, konzernweit gültigen Regelsatz für PMD und SonarQube. Den Regelkatalog stimmte ich mit Vertretern des Konzerns ab und adressierte darin Codequalitätsprobleme sowie sicherheitsrelevante Muster.}
\vspace{0.8mm}

\phtoolline{SonarQube, PMD}

\vspace{1.5mm}
{\color{rulefg}\rule{\linewidth}{0.25pt}\par}
\vspace{2.5mm}


% Project heading
\noindent
{\phprojectname{Pentest}\hfill\phperiod{02/2015 - 03/2015}\par}
{\phprojectmeta{Sogeti NL | Security Tester | Best Practice}%
\hfill {\setlength{\fboxsep}{1.2pt}\colorbox{badgedark}{\textcolor{white}{\fontsize{5.5}{6.0}\selectfont\bfseries Security}}}
\par}

\phcontext{Durchführung eigenständiger Penetrationstests von Webanwendungen und Infrastruktur auf Basis der OWASP-Top-10-Methodik. Dabei kamen Tools wie OWASP ZAP, Burp Suite und Nessus zum Einsatz, um Schwachstellen systematisch zu identifizieren, auszunutzen und zu bewerten. Die Ergebnisse wurden strukturiert dokumentiert, nach Schweregrad priorisiert und mit konkreten technischen Handlungsempfehlungen zur Behebung an die Kunden übergeben.}

\vspace{0.5mm}
\phresptitle{Durchführung von Penetrationstests nach OWASP Top 10}
\phrespdesc{Ich führte eigenständig umfassende Penetrationstests nach der OWASP-Top-10-Methodik durch und setzte dafür OWASP ZAP und Burp Suite für Webanwendungen sowie Nessus und NMap für Infrastruktur- und Netzwerkschwachstellen ein. Damit erstellte ich eine vollständige Sicherheitsbewertung der geprüften Systeme.}
\vspace{0.8mm}
\phresptitle{Reporting und Klassifizierung von Schwachstellen}
\phrespdesc{Ich dokumentierte, klassifizierte und bewertete systematisch alle im Rahmen der Penetrationstests identifizierten Schwachstellen. Ich priorisierte die Befunde nach Schweregrad, beschrieb sie technisch nachvollziehbar in Word und versah sie mit konkreten Handlungsempfehlungen zur gezielten Behebung durch das Entwicklungsteam.}
\vspace{0.8mm}

\phtoolline{OWASP ZAP, Burp Suite, Nmap, Nessus}

\vspace{1.5mm}
{\color{rulefg}\rule{\linewidth}{0.25pt}\par}
\vspace{2.5mm}


% Project heading
\noindent
{\phprojectname{Allegro}\hfill\phperiod{08/2013 - 12/2014}\par}
{\phprojectmeta{Bundesagentur für Arbeit | Software Tester | V-Modell}%
\hfill {\setlength{\fboxsep}{1.2pt}\colorbox{badgeblue}{\textcolor{white}{\fontsize{5.5}{6.0}\selectfont\bfseries AI}}}
\par}

\phcontext{ALLEGRO ist eine zentrale Fachanwendung der Bundesagentur für Arbeit zur Berechnung und Verwaltung von Leistungen der Grundsicherung (Arbeitslosengeld II). Das System basiert auf einer serviceorientierten Architektur und ist in zahlreiche Backend-Systeme wie Zahlungs-, Archiv- und Stammdatensysteme integriert. Es unterstützt Jobcenter bei der Bearbeitung, Berechnung und Auskunft zu Leistungsfällen.}

\vspace{0.5mm}
\phresptitle{Tests zu Sicherheit, Performance, Usability und Barrierefreiheit (NFA)}
\phrespdesc{Ich führte eigenständig ein breites Spektrum nichtfunktionaler Tests durch: Sicherheitstests zur Identifikation von Schwachstellen und Systemstabilität, Usabilitytests zur Bewertung der Benutzerfreundlichkeit sowie Barrierefreiheitstests gemäß BITV-Anforderungen.}
\vspace{0.8mm}
\phresptitle{Weiterentwicklung und Wartung von Testfällen}
\phrespdesc{Ich pflegte, aktualisierte und erweiterte kontinuierlich den vorhandenen Testfallkatalog in SilkCentral entsprechend neuer Anforderungen und Systemänderungen. Ich überarbeitete veraltete Testfälle, ergänzte neue auf Basis von Anforderungsänderungen und verbesserte systematisch die Qualität der gesamten Testfalldokumentation.}
\vspace{0.8mm}
\phresptitle{Barrierefreiheitstests mit Dragon NaturallySpeaking}
\phrespdesc{Ich führte praxisnahe Barrierefreiheitstests unter Einsatz der Spracherkennungssoftware Dragon NaturallySpeaking durch, um die Bedienbarkeit der Anwendung für Menschen mit eingeschränkter Sehstärke sicherzustellen. Die Testergebnisse dokumentierte ich und kommunizierte Verbesserungsempfehlungen an das Entwicklungsteam via SilkCentral.}
\vspace{0.8mm}

\phtoolline{SCTM, SBM, Oracle, Oracle WebLogic, Dragon NaturallySpeaking}

\vspace{1.5mm}
{\color{rulefg}\rule{\linewidth}{0.25pt}\par}
\vspace{2.5mm}


% Project heading
\noindent
{\phprojectname{ProLog}\hfill\phperiod{06/2012 - 05/2013}\par}
{\phprojectmeta{Volkswagen | Software Developer | Wasserfall}%
\hfill {\setlength{\fboxsep}{1.2pt}\colorbox{badgeslate}{\textcolor{white}{\fontsize{5.5}{6.0}\selectfont\bfseries Performance}}}
\par}

\phcontext{ProLog ist eine JEE-basierte Anwendung zur weltweiten Nachverfolgung von Fahrzeugteilen in der Lieferkette des Volkswagen-Konzerns. Das System stellt Transparenz über den aktuellen Standort von Bauteilen bereit – z. B. in LKWs, Flugzeugen oder internationalen Lagern – und unterstützt damit die Just-in-Time-Produktion. Im Projekt erfolgte zudem die Migration der Anwendung von IBM WebSphere auf Apache Tomcat.}

\vspace{0.5mm}
\phresptitle{Weiterentwicklung einer JEE-Logistikanwendung für Volkswagen}
\phrespdesc{Ich analysierte, konzipierte und implementierte neue Funktionalitäten sowie Fehlerbehebungen in einer bestehenden Java-Enterprise-Logistikanwendung. Die Weiterentwicklungen führte ich in enger Zusammenarbeit mit den fachlichen Anforderungsgebern durch und integrierte sie in den ANT-basierten Build-Prozess.}
\vspace{0.8mm}
\phresptitle{Entwicklung und Durchführung von Performancetests mit JMeter}
\phrespdesc{Ich entwickelte eigenständig JMeter-Testpläne zur systematischen Messung und Analyse der Anwendungsperformance unter verschiedenen Lastszenarien. Die Testergebnisse wurden via JavaMelody dargestellt, ich wertete sie aus und nutzte sie als Grundlage für gezielte Performanceoptimierungen an der Anwendung und der Serverinfrastruktur.}
\vspace{0.8mm}
\phresptitle{Serverkonfiguration und automatisierte Deployments mit Shell-Skripten}
\phrespdesc{Ich konfigurierte und verwaltete die Tomcat-Serverumgebung und entwickelte und pflegte Shell-Skripte zur automatisierten Durchführung von Applikations-Deployments. Durch die Automatisierung minimierte ich Deploymentfehler und gestaltete den gesamten Deployment-Prozess reproduzierbar und zuverlässig.}
\vspace{0.8mm}

\phtoolline{Java EE, Tomcat, Oracle, JMeter, JavaMelody, ANT, Subversion, Eclipse, Shell Scripting}

\vspace{1.5mm}
{\color{rulefg}\rule{\linewidth}{0.25pt}\par}
\vspace{2.5mm}


% Project heading
\noindent
{\phprojectname{CP Suite}\hfill\phperiod{04/2011 - 06/2012}\par}
{\phprojectmeta{Corporate Planning | Software Tester | Wasserfall}%
\hfill {\setlength{\fboxsep}{1.2pt}\colorbox{badgeslate}{\textcolor{white}{\fontsize{5.5}{6.0}\selectfont\bfseries Performance}}}
\par}

\phcontext{Corporate Planner (CP-Suite) ist eine Softwareplattform für Corporate Performance Management zur Unterstützung von Controlling, Finanzplanung und Unternehmenssteuerung. Die Lösung ermöglicht integrierte Planung, Budgetierung, Forecasting sowie Reporting und Konsolidierung auf Basis einer zentralen Datenplattform und bietet umfangreiche Analyse- und Reporting-Funktionen für Managemententscheidungen.}

\vspace{0.5mm}
\phresptitle{Entwurf und Durchführung keyword-basierter Tests mit TestComplete}
\phrespdesc{Ich konzipierte, entwickelte und führte keyword-basierte automatisierte Tests mit TestComplete für die CP-Suite-Managementsoftware durch. Der Keyword-Driven-Ansatz ermöglichte mir eine klare Trennung von Testlogik und Testdaten und erleichterte die Wartung und Erweiterung der Testsuiten durch neue Testfälle.}
\vspace{0.8mm}
\phresptitle{Manuelle Tests und Reporting der Befunde in TFS}
\phrespdesc{Ich führte strukturiert manuelle Tests auf Basis definierter Testfälle durch und erfasste, dokumentierte und priorisierte alle festgestellten Defekte im Team Foundation Server (TFS). Die Befundberichte versah ich mit Reproduktionsschritten, Schweregradbewertungen und Screenshots für eine nachvollziehbare Fehlernachverfolgung.}
\vspace{0.8mm}

\phtoolline{TestComplete, TFS}

\vspace{1.5mm}
{\color{rulefg}\rule{\linewidth}{0.25pt}\par}
\vspace{2.5mm}


% Project heading
\noindent
{\phprojectname{Yoom}\hfill\phperiod{07/2010 - 11/2010}\par}
{\phprojectmeta{Yoom | Software Tester | Exploratives Testing, Manuelles Testing und Wasserfall}%
\hfill {\setlength{\fboxsep}{1.2pt}\colorbox{badgegreen}{\textcolor{white}{\fontsize{5.5}{6.0}\selectfont\bfseries Automation}}}
\par}

\phcontext{yoom.de war eine Internetplattform für Wohnungsnachmieten, die Wohnungswechsel mit einer Auktionslogik kombinierte. Ausziehende Mieter konnten Möbel und Inventar zusammen mit ihrem Mietverhältnis einstellen, während Interessenten Gebote abgaben. Die Plattform unterstützte die Vermittlung potenzieller Nachmieter, die anschließend dem Vermieter vorgeschlagen wurden, und digitalisierte damit einen Teil des Wohnungswechsel- und Nachmietprozesses.}

\vspace{0.5mm}
\phresptitle{Manuelle Tests und Reporting der Befunde in Mantis}
\phrespdesc{Ich führte systematisch manuelle Tests der Yoom-Plattform auf Basis definierter Testfälle durch, inklusive funktionaler und explorativer Tests der Kernfunktionalitäten. Alle festgestellten Defekte erfasste ich strukturiert in Mantis, versah sie mit Priorität und Reproduktionsschritten und kommunizierte sie an das Entwicklungsteam.}
\vspace{0.8mm}
\phresptitle{Automatisierte Regressionstests mit Selenium}
\phrespdesc{Ich konzipierte und implementierte automatisierte Regressionstests für die Yoom-Plattform mit Selenium. Ich entwickelte die Tests auf Basis kritischer Nutzerpfade und führte sie regelmäßig aus, um sicherzustellen, dass neue Änderungen keine unbeabsichtigten Rückschritte in der vorhandenen Funktionalität verursachen.}
\vspace{0.8mm}

\phtoolline{Selenium, Mantis}

\vspace{1.5mm}
{\color{rulefg}\rule{\linewidth}{0.25pt}\par}
\vspace{2.5mm}


% Project heading
\noindent
{\phprojectname{Maven Einführung}\hfill\phperiod{05/2009 - 07/2009}\par}
{\phprojectmeta{BITSDONTBYTE | DevOps Engineer | Continuous Integration}%
\hfill {\setlength{\fboxsep}{1.2pt}\colorbox{badgeslate}{\textcolor{white}{\fontsize{5.5}{6.0}\selectfont\bfseries Delivery}}}
\par}

\phcontext{Analyse und Migration bestehender ANT-basierter Build-Konfigurationen auf Apache Maven zur Vereinheitlichung der Build-Prozesse. Erstellung und Pflege von POM-Dateien, Konfiguration von Build-Phasen und Plugins sowie Sicherstellung reproduzierbarer Builds. Zusätzlich Aufbau eines zentralen Maven-Repositories für Bibliotheken und Abhängigkeiten sowie Integration von FishEye zur Verknüpfung von SVN-Commits mit Jira-Tickets und zur verbesserten Nachvollziehbarkeit von Codeänderungen.}

\vspace{0.5mm}
\phresptitle{Migration aller Softwareprojekte von ANT zu Maven}
\phrespdesc{Ich analysierte alle bestehenden ANT-basierten Build-Konfigurationen und migrierte sie schrittweise zu Maven. Ich erstellte POM-Dateien, konfigurierte Build-Phasen und Plugins und stellte sicher, dass alle Projekte nach der Migration korrekt und vollständig gebaut werden konnten.}
\vspace{0.8mm}
\phresptitle{Einführung eines zentralen Repositorys für Bibliotheken und Abhängigkeiten}
\phrespdesc{Ich baute und konfigurierte ein zentrales Maven-Artefakt-Repository zur einheitlichen Verwaltung aller externen und internen Bibliotheken und Abhängigkeiten. Damit löste ich inkonsistente lokale Abhängigkeitsverwaltungen ab und schuf eine reproduzierbare, versionierte und zentral kontrollierte Softwareinfrastruktur für alle Projekte.}
\vspace{0.8mm}
\phresptitle{Einführung von FishEye zur Verknüpfung von SVN-Commits mit Jira-Tickets}
\phrespdesc{Ich installierte, konfigurierte und integrierte das FishEye-Plugin in die bestehende Jira-Instanz, um eine direkte Verknüpfung zwischen SVN-Code-Commits und den zugehörigen Jira-Tickets herzustellen. Damit verbesserte ich die Nachvollziehbarkeit von Codeänderungen und deren Bezug zu Anforderungen deutlich.}
\vspace{0.8mm}

\phtoolline{ANT, Eclipse, Maven, FishEye, Jira, Subversion}

\vspace{1.5mm}
{\color{rulefg}\rule{\linewidth}{0.25pt}\par}
\vspace{2.5mm}


% Project heading
\noindent
{\phprojectname{MISSLight}\hfill\phperiod{12/2005 - 12/2006}\par}
{\phprojectmeta{Volkswagen | Software Developer | inkrementelles Modell}%
\hfill {\setlength{\fboxsep}{1.2pt}\colorbox{badgeslate}{\textcolor{white}{\fontsize{5.5}{6.0}\selectfont\bfseries Performance}}}
\par}

\phcontext{Entwicklung und Betrieb einer JEE-basierten Anwendung zur konzernweiten Verwaltung von Fahrzeugteilen. In dem System wurden sämtliche Bauteile des Volkswagen-Konzerns mit ihren Eigenschaften, Varianten und technischen Merkmalen zentral gepflegt. Die Anwendung unterstützte die strukturierte Verwaltung und Suche nach Material- und Stammdaten und stellte diese Informationen verschiedenen Entwicklungs- und Produktionssystemen zur Verfügung. Ein solches Teilemanagement dient dazu, Bauteile und ihre Eigenschaften zentral zu verwalten und effizient nutzbar zu machen.}

\vspace{0.5mm}
\phresptitle{Schnittstellenerweiterungen der Teileverwaltungssoftware}
\phrespdesc{Ich analysierte bestehende Schnittstellen der Teileverwaltungssoftware und implementierte Erweiterungen zur Anbindung neuer Systeme und zur Anpassung an geänderte Anforderungen. Dabei berücksichtigte ich die bestehende Systemarchitektur und testete die Änderungen sorgfältig, um die Stabilität der Produktivumgebung sicherzustellen.}
\vspace{0.8mm}
\phresptitle{Performanceoptimierungen der JEE-Anwendung}
\phrespdesc{Ich identifizierte und behob Performanceengpässe in der JEE-Anwendung. Dazu analysierte und optimierte ich Datenbankabfragen, überarbeitete ineffiziente Codeabschnitte und passte serverseitige Konfigurationsparameter an. Die erzielten Verbesserungen quantifizierte und dokumentierte ich durch Messungen vor und nach der Optimierung.}
\vspace{0.8mm}
\phresptitle{Serverkonfiguration und Deployments auf WebSphere}
\phrespdesc{Ich konfigurierte den WebSphere-Applikationsserver und plante und führte Applikations-Deployments eigenverantwortlich in der Produktivumgebung durch. Durch strukturierte Deployment-Prozesse und sorgfältige Vorbereitung minimierte ich Ausfallzeiten und stellte die Verfügbarkeit des Systems für den Produktionsbetrieb sicher.}
\vspace{0.8mm}

\phtoolline{Java EE, WebSphere, Oracle, Excel, Eclipse, Toad, ANT, Subversion}

\vspace{1.5mm}
{\color{rulefg}\rule{\linewidth}{0.25pt}\par}
\vspace{2.5mm}


% Project heading
\noindent
{\phprojectname{Umwelt-Informationssystem}\hfill\phperiod{10/2005 - 12/2005}\par}
{\phprojectmeta{Volkswagen | Software Developer | Wasserfall}%
\hfill {\setlength{\fboxsep}{1.2pt}\colorbox{badgeslate}{\textcolor{white}{\fontsize{5.5}{6.0}\selectfont\bfseries Delivery}}}
\par}

\phcontext{Entwicklung und Betrieb eines Systems zur Bewertung der Umweltverträglichkeit und Recyclingfähigkeit von Fahrzeugen. Auf Basis der Material- und Bauteildaten einzelner Fahrzeugkomponenten wurde der jeweilige Recyclinganteil berechnet und aggregiert. Dadurch konnte die Gesamt-Recyclingquote eines Fahrzeugs ermittelt werden und als Grundlage für Umweltbewertungen und Nachhaltigkeitsanforderungen im Konzern dienen.}

\vspace{0.5mm}
\phresptitle{Entwicklung eines Fat-Clients mit Java Swing}
\phrespdesc{Ich konzipierte und implementierte ein Fat-Client-Frontend auf Basis von Java Swing in Eclipse für das Umwelt-Informationssystem. Das Frontend ermöglichte die komfortable Erfassung und Analyse von Fahrzeugkomponentendaten zur Bewertung der Recyclebarkeit und wurde gezielt im Hinblick auf Benutzerfreundlichkeit optimiert. Das Backend wurde mit Spring und Hibernate (Java EE) implementiert.}
\vspace{0.8mm}
\phresptitle{Entwicklung eines Frameworks für wiederverwendbare UI-Komponenten}
\phrespdesc{Ich entwarf und implementierte ein wiederverwendbares Framework für standardisierte UI-Komponenten auf Basis von Java Swing. Das Framework stellte häufig benötigte Eingabe-, Ausgabe- und Navigationskomponenten bereit, förderte eine einheitliche Benutzeroberfläche und reduzierte den Entwicklungsaufwand für neue Systemfunktionalitäten erheblich.}
\vspace{0.8mm}

\phtoolline{Java EE, Java Swing, Oracle, WebSphere, Toad, ANT, Subversion, Eclipse}

\vspace{1.5mm}
{\color{rulefg}\rule{\linewidth}{0.25pt}\par}
\vspace{2.5mm}


% Project heading
\noindent
{\phprojectname{KWG}\hfill\phperiod{08/2005 - 10/2005}\par}
{\phprojectmeta{Deutsche Bundesbank | Software Developer | inkrementelles Modell}%
\hfill {\setlength{\fboxsep}{1.2pt}\colorbox{badgegreen}{\textcolor{white}{\fontsize{5.5}{6.0}\selectfont\bfseries Automation}}}
\par}

\phcontext{Entwicklung einer webbasierten Anwendung zur automatisierten Bewertung der Kreditwürdigkeit von Antragstellern auf Grundlage regulatorischer Vorgaben zur Kreditwürdigkeitsprüfung. Die serverseitige Implementierung basierte auf Apache Struts. Das System analysierte relevante Finanz- und Risikodaten und unterstützte eine strukturierte Entscheidung über die Kreditfähigkeit von Antragstellern.}

\vspace{0.5mm}
\phresptitle{Erweiterung und Wartung von JEE-Anwendungen der Deutschen Bundesbank}
\phrespdesc{Ich analysierte, erweiterte und behob Fehler in bestehenden Java-Enterprise-Anwendungen der Deutschen Bundesbank zur Kreditwürdigkeitsbewertung. Ich implementierte neue Anforderungen, behob gemeldete Fehler und stellte den stabilen und zuverlässigen Betrieb der geschäftskritischen Anwendung im produktiven Umfeld sicher. Die Serverlogik wurde mit Java EE und Struts implementiert.}
\vspace{0.8mm}

\phtoolline{Java EE, Oracle WebLogic, Oracle, Struts, Eclipse}

\vspace{1.5mm}
{\color{rulefg}\rule{\linewidth}{0.25pt}\par}
\vspace{2.5mm}


% Project heading
\noindent
{\phprojectname{Dozent}\hfill\phperiod{10/2004 - 07/2005}\par}
{\phprojectmeta{Volkswagen | Software Developer | Agile Vorgehensweise}%
\hfill {\setlength{\fboxsep}{1.2pt}\colorbox{badgeslate}{\textcolor{white}{\fontsize{5.5}{6.0}\selectfont\bfseries Delivery}}}
\par}

\phcontext{DOZENT war eine JEE-basierte Enterprise-Anwendung zur Verwaltung von Fahrzeugteilen im Volkswagen-Konzern. Das System diente der zentralen Pflege und Verwaltung von Bauteilen sowie deren Eigenschaften und Varianten. Es unterstützte Fachbereiche dabei, Material- und Produktinformationen strukturiert zu erfassen, zu pflegen und für Entwicklungs- und Produktionsprozesse bereitzustellen.}

\vspace{0.5mm}
\phresptitle{Implementierung von Erweiterungen in der Dozent-Anwendung}
\phrespdesc{Ich analysierte bestehende Anforderungen und implementierte neue Funktionalitäten in der Dozent-Anwendung auf Basis von Java EE. Dabei berücksichtigte ich die bestehende Systemarchitektur und Datenmodelle, testete die Erweiterungen sorgfältig und spielte sie in die produktive Umgebung ein.}
\vspace{0.8mm}
\phresptitle{Durchführung von Bugfixes in der Verwaltungssoftware}
\phrespdesc{Ich identifizierte, analysierte und behob gemeldete Fehler in der Verwaltungssoftware für Automobilteile. Dabei reproduzierte ich die Fehler systematisch, führte eine Ursachenanalyse im bestehenden Codebestand durch und implementierte und verifizierte geeignete Korrekturen, um den stabilen Betrieb der Anwendung sicherzustellen.}
\vspace{0.8mm}

\phtoolline{Java EE, OC4J, Oracle, Eclipse, ANT, CVS}

\vspace{1.5mm}
{\color{rulefg}\rule{\linewidth}{0.25pt}\par}
\vspace{2.5mm}


% Project heading
\noindent
{\phprojectname{ESV}\hfill\phperiod{02/2004 - 07/2004}\par}
{\phprojectmeta{Volkswagen | Software Developer | Agile Vorgehensweise}%
\hfill {\setlength{\fboxsep}{1.2pt}\colorbox{badgeslate}{\textcolor{white}{\fontsize{5.5}{6.0}\selectfont\bfseries Delivery}}}
\par}

\phcontext{Entwicklung einer Software zur Unterstützung und Überwachung manueller Montageprozesse in der Produktion. Mitarbeitende erhielten Arbeitskörbe mit Bauteilen und wurden über einen großen Bildschirm mit Ampelanzeige über ihre aktuelle Arbeitsgeschwindigkeit informiert. Die Anwendung visualisierte den Bearbeitungsfortschritt in Echtzeit und zeigte über eine Grün-Gelb-Rot-Logik an, ob der Arbeitsrhythmus den Produktionsvorgaben entspricht.}

\vspace{0.5mm}
\phresptitle{Entwicklung eines Prototyps auf neuer technologischer Basis}
\phrespdesc{Ich konzipierte und implementierte einen funktionsfähigen Prototyp für das Produktionssteuerungssystem ESV auf Basis moderner Java-EE-Technologien EJB auf einem JBoss-Applicationserver. Damit demonstrierte ich die technische Machbarkeit des neuen Architekturansatzes und legte die Grundlage für die weitere Produktentwicklung im Volkswagen-Konzern.}
\vspace{0.8mm}

\phtoolline{Eclipse, JBoss, Tomcat, Java EE (EJB)}

\vspace{1.5mm}


\end{document}